\documentclass[12pt,a4paper]{article}

\usepackage[utf8]{inputenc}
\usepackage[T2A]{fontenc}
\usepackage[ukrainian]{babel}
\usepackage{amsmath,amssymb}
\usepackage{array}
\usepackage{booktabs}
\usepackage{tabularx}
\usepackage{longtable}
\usepackage{geometry}
\usepackage{hyperref}

\geometry{left=15mm,right=15mm,top=15mm,bottom=15mm}
\pagestyle{empty}
\setlength{\parindent}{0pt}
\setlength{\parskip}{6pt}
\renewcommand{\arraystretch}{1.25}

\begin{document}

\section*{Порівняльний аналіз методів виділення універсальних ознак
у моделях простору станів}

Моделі простору станів, або State Space Models, далі SSM, істотно
відрізняються від класичних згорткових мереж і трансформерів тим, що
їхні внутрішні представлення формуються не через локальні згорткові
фільтри або явні матриці уваги, а через латентну динаміку стану.

У лінійному випадку така модель задається рівняннями
\[
    x'(t)=Ax(t)+Bu(t), \qquad y(t)=Cx(t)+Du(t),
\]
де \(u(t)\) є вхідним сигналом, \(x(t)\) є латентним станом, а
\(y(t)\) є виходом моделі. Після дискретизації така система може
розглядатися як рекурентна модель або як згортка з імпульсною
характеристикою
\[
    K_k = C A^k B.
\]

Саме ця подвійність створює головну відмінність SSM від CNN та ViT:
ознаки доцільно шукати не лише у прихованих активаціях, а також у
ядрах, модах, власних значеннях, полюсах, імпульсних характеристиках
і, для selective SSM, у вхідно-залежних маршрутах передавання
інформації.

\section*{1. Порівняння сімейств SSM}

Різні сімейства SSM мають різну структуру латентної динаміки, тому
природний рівень аналізу ознак для них також відрізняється. Для
лінійних часово-інваріантних моделей основними об'єктами аналізу є
ядра, моди та частотні характеристики, тоді як для selective SSM
важливішими є вхідно-залежні патерни маршрутизації інформації.

\begin{longtable}{p{2.6cm} p{3.6cm} p{3.4cm} p{3.2cm} p{3.4cm}}
\toprule
\textbf{Сімейство}
&
\textbf{Структурна основа}
&
\textbf{Кандидат на ознаку}
&
\textbf{Інтерпретаційний зміст}
&
\textbf{Основне обмеження}
\\
\midrule
\endfirsthead

\toprule
\textbf{Сімейство}
&
\textbf{Структурна основа}
&
\textbf{Кандидат на ознаку}
&
\textbf{Інтерпретаційний зміст}
&
\textbf{Основне обмеження}
\\
\midrule
\endhead

HiPPO та ранні LTI SSM
&
Лінійна часово-інваріантна система, що апроксимує історію сигналу
через ортогональні базисні проєкції.
&
Коефіцієнт базисної проєкції або згладжена історична компонента.
&
Часовий масштаб, тренд, довготривала пам'ять.
&
Коефіцієнти часто є абстрактними і не мають прямої семантичної
інтерпретації.
\\

S4
&
Структурована SSM з матрицею стану типу diagonal plus low-rank.
&
Ядро згортки, полюс, власна мода, полюсно-резидуальна пара.
&
Довготривала залежність, лаг, частотний фільтр.
&
Сирі координати стану не є стабільними семантичними одиницями.
\\

DSS та S4D
&
Діагональна або спрощена діагональна параметризація SSM.
&
Окрема діагональна мода або частотна смуга.
&
Чистіший часовий або спектральний детектор.
&
Навіть діагональні моди можуть змішувати різні семантичні фактори.
\\

S5
&
MIMO-SSM, тобто один багатовхідний і багатовихідний простір станів.
&
Спільна мода, що діє на кілька вхідних і вихідних каналів.
&
Міжканальний часовий мотив.
&
MIMO-змішування ускладнює відповідність між однією координатою і
одним поняттям.
\\

LRU
&
Діагональна лінійна рекурентна модель з нормалізованою динамікою.
&
Рекурентний канал або стабільна власна мода.
&
Довготривалий лічильник, структурна пам'ять, акумуляція історії.
&
Потрібне зіставлення каналів між різними запусками навчання.
\\

Mamba
&
Selective SSM, у якій параметри динаміки залежать від поточного входу.
&
Вхідно-залежний маршрут, gate-патерн або dependency motif.
&
Контентне маршрутизування, вибіркове збереження або забування
інформації.
&
Один і той самий канал може виконувати різні ролі в різних контекстах.
\\

ViM
&
Візуальна модель на основі bidirectional Mamba blocks.
&
Просторова карта після обернення послідовності патчів у 2D-координати.
&
Границі, частини об'єкта, глобальна форма.
&
Порядок сканування може створювати артефакти.
\\

VMamba
&
Візуальна SSM з 2D selective scan.
&
Маршрутно-залежна просторова ознака або dependency pattern.
&
Просторовий контекст, глобальна взаємодія областей зображення.
&
Потрібна перевірка ознак для різних напрямів сканування.
\\
\bottomrule
\end{longtable}

\section*{2. Критерії коректного виділення універсальної ознаки}

Ознаку в SSM не можна вважати універсальною лише тому, що вона дає
високу активацію або добре працює в одному запуску навчання. Коректна
перевірка повинна враховувати семантичну узгодженість, стабільність
між незалежно навченими моделями, каузальний вплив на результат,
переносимість на інші задачі та, для selective SSM, сталість
вхідно-залежних маршрутів.

Формально таку перевірку можна подати як набір критеріїв
\[
    \Phi = (A,S,C,T,D),
\]
де \(A\) позначає узгодженість з інтерпретованим поняттям,
\(S\) позначає стабільність між різними seed-моделями,
\(C\) позначає каузальну важливість,
\(T\) позначає переносимість на інші задачі,
а \(D\) позначає стабільність залежностей для selective SSM.

\begin{center}
\begin{tabularx}{\textwidth}{p{2cm} X X}
\toprule
\textbf{Критерій}
&
\textbf{Зміст}
&
\textbf{Практична перевірка}
\\
\midrule

\(A\)
&
Alignment: ознака має відповідати інтерпретованому поняттю.
&
IoU з масками понять, span-IoU для тексту або часових рядів,
узгодження з частотними областями для аудіо.
\\

\(S\)
&
Stability: ознака має повторюватися у незалежно навчених моделях.
&
Зіставлення каналів, мод або підпросторів між різними seed-моделями.
\\

\(C\)
&
Causal importance: ознака повинна впливати на поведінку моделі.
&
Ablation, intervention, deletion AUC, падіння якості після вимкнення
ознаки.
\\

\(T\)
&
Transfer utility: ознака повинна бути корисною на іншій задачі.
&
Лінійний probe або простий класифікатор на заморожених ознаках.
\\

\(D\)
&
Dependency consistency: для selective SSM має зберігатися подібний
патерн залежностей.
&
Порівняння induced attention, relevance maps або gate-conditioned paths.
\\
\bottomrule
\end{tabularx}
\end{center}

Критерій \(D\) застосовується лише до selective SSM, зокрема Mamba,
ViM та VMamba. Для лінійних часово-інваріантних SSM достатньо критеріїв
\(A,S,C,T\), але бажано додатково аналізувати частотну стабільність
мод.

\section*{3. Чому сирі координати стану не є надійними ознаками}

Для SSM існує проблема неідентифікованості стану. Якщо \(T\) є
оборотною матрицею, то заміна
\[
    \tilde{x}(t)=Tx(t)
\]
дає еквівалентну систему
\[
    \tilde{A}=TAT^{-1}, \qquad
    \tilde{B}=TB, \qquad
    \tilde{C}=CT^{-1}.
\]
Вхідно-вихідна поведінка при цьому не змінюється. Отже, окрема
координата \(x_j(t)\) не має гарантованого семантичного змісту.

Через це в SSM безпечніше аналізувати не сирі координати стану, а
інваріантні або краще визначені об'єкти: ядра \(K_k\), власні моди,
полюси, імпульсні характеристики, частотні відповіді, підпростори
після вирівнювання між seed-моделями та gate-conditioned dependency
patterns.

Для діагоналізованої SSM внесок окремої моди можна записати як
\[
    K_k^{(m)} = \alpha_m \lambda_m^k \beta_m,
    \qquad
    K_k = \sum_m K_k^{(m)}.
\]
Тут \(\lambda_m\) є власним значенням або дискретним полюсом,
а \(\alpha_m\) та \(\beta_m\) описують проєкції на вихідний і вхідний
простори. Саме така форма є коректною для аналізу мод, оскільки вона
не зводить інтерпретацію до довільно вибраної координати стану.

\section*{4. Порівняння методів інтерпретації}

Методи інтерпретації, розроблені для CNN та ViT, не можна механічно
переносити на SSM. Частина з них працює після адаптації, частина
потребує суттєвої зміни об'єкта аналізу, а деякі методи мають сенс
лише для лінійних часово-інваріантних SSM.

\begin{longtable}{p{3.2cm} p{2.7cm} p{3cm} p{3cm} p{4.2cm}}
\toprule
\textbf{Метод}
&
\textbf{CNN / ViT}
&
\textbf{LTI SSM}
&
\textbf{Selective SSM}
&
\textbf{Необхідна модифікація}
\\
\midrule
\endfirsthead

\toprule
\textbf{Метод}
&
\textbf{CNN / ViT}
&
\textbf{LTI SSM}
&
\textbf{Selective SSM}
&
\textbf{Необхідна модифікація}
\\
\midrule
\endhead

Network Dissection
&
Природне застосування до карт активацій.
&
Можливе, але слід працювати з часовими або частотними масками.
&
Можливе для ViM та VMamba після обернення scan-order у 2D.
&
Піксельні маски треба замінити на часові, частотні, span-маски або
inverse-scan spatial masks.
\\

Activation Maximization
&
Добре працює для візуалізації фільтрів і нейронів.
&
Менш надійне, бо синтетичний сигнал може створювати артефакти.
&
Обмежено корисне через контекстну залежність gating.
&
Потрібні регуляризатори гладкості, спектральні обмеження і перевірка
на природних прикладах.
\\

Аналітичний аналіз ядра
&
Прямого аналога немає.
&
Найприродніший метод для S4, DSS, S4D, S5, LRU.
&
Малопридатний, бо параметри залежать від входу.
&
Будуються impulse response, pole map, transfer function, Bode magnitude.
\\

Exemplar Retrieval
&
Вибір природних прикладів з максимальною активацією.
&
Потрібно відбирати приклади за домінуванням мод або ядра.
&
Потрібно відбирати приклади за dependency maps або relevance maps.
&
Кластеризацію слід виконувати не в сирому просторі станів, а в просторі
узгоджених ознак.
\\

Circuit Analysis
&
Аналіз локальних нейронних схем, фільтрів і зв'язків.
&
Аналіз композиції мод, полюсів і міжшарових ядер.
&
Аналіз gate-conditioned routes.
&
Потрібно переходити від статичних нейронів до динамічних маршрутів.
\\

Attribution / LRP
&
Стандартні карти релевантності.
&
Можливе з обережністю, але слід враховувати рекурентну динаміку.
&
Особливо важливе для Mamba-подібних моделей.
&
Потрібні SSM-специфічні або induced-attention підходи.
\\

Linear Probes
&
Простий і сильний базовий метод.
&
Добре працює для pooled states, modes або kernel features.
&
Добре працює, якщо probe враховує контекстні залежності.
&
Бажано використовувати регуляризацію, щоб уникнути змішування ознак.
\\

Контроль порядку сканування
&
Зазвичай не потрібний.
&
Не застосовується до звичайних 1D SSM.
&
Обов'язковий для ViM та VMamba.
&
Ознака має перевірятися при різних scan routes.
\\
\bottomrule
\end{longtable}

\section*{5. Структурне порівняння CNN, ViT, LTI SSM та selective SSM}

Структурне порівняння потрібне для того, щоб показати, на якому
рівні кожна архітектура природно формує інтерпретовані ознаки.
У CNN такими одиницями часто є просторові фільтри, у ViT —
token-to-token взаємодії, у LTI SSM — моди та ядра, а в selective SSM —
динамічні маршрути, що залежать від вхідного сигналу.

\begin{longtable}{p{3cm} p{3.2cm} p{3.2cm} p{3.2cm} p{3.2cm}}
\toprule
\textbf{Вимір порівняння}
&
\textbf{CNN}
&
\textbf{ViT}
&
\textbf{LTI SSM}
&
\textbf{Selective SSM}
\\
\midrule
\endfirsthead

\toprule
\textbf{Вимір порівняння}
&
\textbf{CNN}
&
\textbf{ViT}
&
\textbf{LTI SSM}
&
\textbf{Selective SSM}
\\
\midrule
\endhead

Локальність ознак
&
Просторово локальна, receptive field зростає з глибиною.
&
Глобальна через self-attention.
&
Історія інтегрується через ядро.
&
Історія інтегрується вибірково залежно від вхідного сигналу.
\\

Базис ознак
&
Фільтри згортки.
&
Key-query-value взаємодії.
&
Моди, полюси, ядра, власні значення.
&
Вхідно-залежні маршрути та gating patterns.
\\

Ідентифікованість стану
&
Відносно висока: фільтр має стабільну просторову роль.
&
Часткова: attention heads можуть мати ролі, але вони не завжди сталі.
&
Низька для сирих координат, вища для мод і ядер.
&
Низька для каналів, бо семантика залежить від контексту.
\\

Частотний характер
&
Збалансований, залежить від даних і архітектури.
&
Може залежати від позиційного кодування і навчання.
&
Часто схильний до низькочастотних і довготривалих залежностей.
&
Адаптивний, визначається selective gates.
\\

Семантика глибини
&
Перехід від країв і текстур до частин та об'єктів.
&
Менш жорстка ієрархія, роль шару залежить від архітектури.
&
Глибина не обов'язково відповідає просторовій ієрархії.
&
Глибина радше відображає складність маршрутизації інформації.
\\

Найкращий рівень аналізу
&
Середні шари та карти активацій.
&
Attention heads, token embeddings, MLP features.
&
Kernel modes, pole-residue pairs, impulse responses.
&
Dependency motifs, gate-conditioned paths, induced attention.
\\
\bottomrule
\end{longtable}

\section*{6. Порівняння обмежень}

Обмеження SSM-UFE мають переважно архітектурну природу. Вони не
зводяться лише до нестачі інструментів інтерпретації, оскільки
випливають із самої будови SSM: латентний стан не є однозначно
ідентифікованим, частотна поведінка часто зміщена до довготривалих
залежностей, а selective SSM змінюють свої маршрути залежно від входу.

\begin{longtable}{p{3.5cm} p{3cm} p{2.7cm} p{2.7cm} p{4.2cm}}
\toprule
\textbf{Обмеження}
&
\textbf{Кого стосується}
&
\textbf{Серйозність для CNN / ViT}
&
\textbf{Серйозність для SSM}
&
\textbf{Коректна стратегія}
\\
\midrule
\endfirsthead

\toprule
\textbf{Обмеження}
&
\textbf{Кого стосується}
&
\textbf{Серйозність для CNN / ViT}
&
\textbf{Серйозність для SSM}
&
\textbf{Коректна стратегія}
\\
\midrule
\endhead

Неідентифікованість стану
&
Усі SSM.
&
Низька.
&
Висока.
&
Не інтерпретувати сирі координати стану; аналізувати ядра, моди,
полюси та вирівняні підпростори.
\\

Низькочастотний зсув
&
Насамперед LTI SSM.
&
Низька або середня.
&
Середня або висока.
&
Додавати частотні ablation-тести і перевіряти високочастотні ознаки.
\\

Контекстна залежність семантики
&
Selective SSM.
&
Середня для attention-моделей.
&
Висока.
&
Формулювати універсальність не на рівні каналу, а на рівні
dependency motif.
\\

Відсутність просторової ієрархії
&
Усі SSM.
&
Низька.
&
Середня.
&
Вибирати шари через kernel analysis, а не через CNN-евристику
``чим глибше, тим семантичніше''.
\\

Артефакти порядку сканування
&
ViM та VMamba.
&
Низька.
&
Висока.
&
Перевіряти різні scan routes і повертати ознаки у 2D-простір.
\\

Масштабування навчання ознак
&
Великі SSM.
&
Середня.
&
Середня або висока.
&
Порівнювати кілька ширин моделі, seed-моделей і параметризацій.
\\
\bottomrule
\end{longtable}

\section*{7. Рекомендований протокол порівняльного експерименту}

Для коректного порівняння CNN, ViT, LTI SSM та selective SSM потрібно
оцінювати не лише якість на основній задачі, а й стабільність,
інтерпретованість, каузальну значущість і переносимість знайдених
ознак. Тому експериментальний протокол має включати кілька типів
даних і кілька незалежних метрик.

\begin{longtable}{p{3cm} p{3cm} p{3.2cm} p{3.4cm} p{3.6cm}}
\toprule
\textbf{Напрям}
&
\textbf{Моделі}
&
\textbf{Дані}
&
\textbf{Ознаки}
&
\textbf{Метрики}
\\
\midrule
\endfirsthead

\toprule
\textbf{Напрям}
&
\textbf{Моделі}
&
\textbf{Дані}
&
\textbf{Ознаки}
&
\textbf{Метрики}
\\
\midrule
\endhead

Зорові ознаки
&
CNN, ViT, ViM, VMamba.
&
ImageNet, Broden, COCO, ADE20K.
&
Краї, текстури, частини об'єктів, scan-direction motifs.
&
IoU, deletion AUC, cross-seed matching, transfer accuracy.
\\

Довгі послідовності
&
S4, S5, DSS, LRU, Mamba, Transformer.
&
Long Range Arena, Pathfinder, Path-X.
&
Довготривалі залежності, лічильники, дужкові структури, path continuity.
&
Probe accuracy, ablation drop, length extrapolation.
\\

Аудіо та часові ряди
&
S4D, DSS, S5, LRU, Mamba.
&
Speech Commands, часові ряди.
&
Onsets, pitch bands, keyword spans, spectral modes.
&
Span-IoU, spectral selectivity, perturbation AUC.
\\

Перенесення ознак
&
Будь-які заморожені SSM, CNN або ViT.
&
Малі downstream-задачі.
&
Канали, моди, pooled states, dependency motifs.
&
Linear accuracy, sample efficiency, transfer gap.
\\
\bottomrule
\end{longtable}

\section*{8. Узагальнена порівняльна оцінка}

Узагальнена оцінка показує, що різні архітектури мають різний
природний рівень інтерпретації. Для CNN таким рівнем є карти
активацій і згорткові фільтри, для ViT — token-представлення та
механізми уваги, для лінійних SSM — ядра, моди й полюси, а для
selective SSM — вхідно-залежні маршрути передавання інформації.

\begin{center}
\begin{tabularx}{\textwidth}{p{3.2cm} X X X}
\toprule
\textbf{Архітектура}
&
\textbf{Сильна сторона для UFE}
&
\textbf{Слабка сторона для UFE}
&
\textbf{Найбільш коректний об'єкт аналізу}
\\
\midrule

CNN
&
Стабільні локальні фільтри і природна просторова ієрархія.
&
Менш придатні для дуже довгих залежностей.
&
Карти активацій, фільтри, середні шари.
\\

ViT
&
Глобальні token-to-token взаємодії.
&
Attention не завжди є повним поясненням причинності.
&
Attention heads, token embeddings, attribution maps.
\\

LTI SSM
&
Аналітично доступні ядра, полюси, моди і частотні характеристики.
&
Сирі координати стану неідентифіковані.
&
Impulse response, transfer function, pole-residue pairs.
\\

Selective SSM
&
Вхідно-залежне маршрутизування інформації і лінійна складність за
довжиною послідовності.
&
Семантика каналів змінюється від контексту до контексту.
&
Gate-conditioned dependency motifs, induced attention, relevance maps.
\\

Vision SSM
&
Поєднання глобального контексту з ефективним скануванням зображення.
&
Можливі артефакти порядку сканування.
&
Inverse-scan spatial maps і route-consistent features.
\\
\bottomrule
\end{tabularx}
\end{center}

\section*{Висновки}

Універсальні ознаки в SSM не слід ототожнювати із сирими координатами
латентного стану. На відміну від CNN, де окремий фільтр часто можна
трактувати як відносно стабільний просторовий детектор, у SSM стан
визначений з точністю до зміни базису. Тому для лінійних SSM найбільш
надійними об'єктами аналізу є ядра, власні моди, імпульсні
характеристики, полюси та частотні відповіді.

Для selective SSM, зокрема Mamba-подібних моделей, універсальність
доцільно формулювати не як сталість одного каналу, а як повторюваність
функціонального маршруту: gate-conditioned dependency motif, induced
attention pattern або relevance map. Такий підхід є обережнішим, але
науково коректнішим, оскільки враховує вхідно-залежну природу
selective state space dynamics.

Порівняння з CNN та ViT показує, що частина класичних методів
інтерпретації переноситься на SSM лише після модифікації. Network
Dissection, exemplar retrieval, linear probes та attribution-методи
можуть бути корисними, але потребують адаптації до часових,
частотних або scan-inverted просторових масок. Натомість аналіз
імпульсних характеристик, полюсів і частотних відповідей є специфічною
перевагою LTI SSM і не має прямого аналога у класичних CNN або ViT.

Отже, найбільш обґрунтованим є не твердження про абсолютну
універсальність ознак у всіх архітектурах, а твердження про
архітектурно-залежні універсальні функціональні мотиви, які
стабільно виникають у межах певного сімейства SSM, узгоджуються з
інтерпретованим поняттям, мають каузальний вплив на результат моделі
і зберігають корисність при перенесенні на інші задачі.



\end{document}